% UNIFIED 2026-05-13
% =====================================================================
% Preliminaries -- Geometric Measure Theory
% Agent 2 deliverable. To be merged by Agent 17 into
%     Manuscript/01_preliminaries_draft.tex.
%
% Notation: follows Manuscript/00_notation_register.md verbatim.
% Macros assumed loaded (per register, §15): \R, \Fr, \Om, \eps,
%     \rstar, \rhad, \rhodelta, \zbar, \X, \dX, \Bstar, \Erho, \Frho,
%     \Vrho, \nurho, \deltaT, \deltaFK.
%
% SOURCE MAP for this section (each item is one of:
%   PROVED HERE / CITED FROM <ref> / ADAPTED FROM <local file>):
%
%   §A.1  Lebesgue/Hausdorff measure, perimeter, reduced boundary
%         and normal $\nurho$
%             CITED FROM Maggi, "Sets of Finite Perimeter and
%             Geometric Variational Problems" (Cambridge 2012),
%             Def. 12.1, Def. 15.1, Thm. 15.9.
%             CITED FROM Ambrosio--Fusco--Pallara (AFP),
%             "Functions of Bounded Variation and Free Discontinuity
%             Problems" (Oxford 2000), Def. 3.1, Thm. 3.59.
%
%   §A.2  Lower semicontinuity of perimeter under $L^1_{\rm loc}$ conv.
%             CITED FROM Maggi 2012, Prop. 12.15.
%             CITED FROM AFP 2000, Thm. 3.7.
%
%   §A.3  BV coarea formula (statement + integrated form)
%             CITED FROM Maggi 2012, Thm. 13.1 (BV coarea),
%             and Thm. 18.1 (Sobolev coarea).
%             CITED FROM AFP 2000, Thm. 3.40.
%
%   §A.4  Sobolev coarea applied to the torsion function $u_\Om$ and
%         Sobolev--Sard a.e. regular-level statement giving
%         $\Vrho=-t'(\rho)/|\nabla u|$ a.e.
%             CITED FROM Figalli--Maggi, "On the shape of liquid drops
%             and crystals in the small mass regime", Arch. Ration.
%             Mech. Anal. 201 (2011), App. A, Thm. A.1.
%             ADAPTED FROM Plan 2/WIP/WIP_LevelSetIdentity.tex
%             §2 (coarea applicability) and §4.
%
%   §A.5  Divergence theorem on sets of finite perimeter
%         (Gauss--Green for $C^1_c$ fields and for $C^1$ fields
%         bounded on $E$ together with $\Div X$ in $L^1(E)$).
%             CITED FROM Maggi 2012, Thm. 15.5 (Gauss--Green
%             for sets of finite perimeter) and Thm. 16.3 (its
%             $C^1$/$L^1$ form via approximation).
%
%   §A.6  Lipschitz truncation flux identity for the torsion function
%         (replacing the boundary trace of $\nabla u$ on
%         $\partial^*\Erho$ by the BV coarea+Lipschitz cutoff).
%             ADAPTED FROM Plan 2/WIP/WIP_LevelSetIdentity.tex,
%             Lem. "Flux = volume" (lem:flux) and Rem. rem:divergence.
%
%   §A.7  Truncation lemma for singular vector fields of the type
%         $X(x)=w(|x-z|)(x-z)/|x-z|$ with $w$ bounded.
%             PROVED HERE.
%             Used by Agent 11 in the oriented radial trace lemma
%             (notation register §9: $T_\rho$). The proof is the
%             rigorous form of the truncation argument written
%             informally in Plan 2/WIP/WIP_WeightedMetricTrace.tex
%             (proof of lem:radial-slicing, lines 182--195) and is
%             load-bearing.
%
%   §A.8  Reshetnyak lower semicontinuity for linear-growth
%         integrands on sets of finite perimeter; passage of the
%         finite-perimeter first variation from smooth flows.
%             CITED FROM AFP 2000, Thm. 2.38 (Reshetnyak lower
%             semicontinuity).
%             ADAPTED FROM Plan 2/gmt-inputs-for-metric-closure.md §5
%             (the four-step approximation procedure).
%
%   §A.9  Measurable selection of levelwise centres
%             PROVED HERE (via Aumann/Kuratowski--Ryll-Nardzewski
%             selection).
%             CITED FROM Castaing--Valadier, "Convex Analysis and
%             Measurable Multifunctions" (Springer LNM 580, 1977),
%             Thm. III.6 and Thm. III.8;
%             alternatively Federer, "Geometric Measure Theory"
%             (Springer 1969), §2.13.
%             ADAPTED FROM Plan 2/WIP/WIP_WeightedMetricTrace.tex
%             (proof of lem:FJ-package, last paragraph).
%
% =====================================================================

\subsection{Preliminaries: geometric measure theory}\label{sec:prelim-gmt}

This section fixes the finite-perimeter, slicing, and divergence-theorem
tools used throughout the manuscript.  All notation matches
\S\ref{sec:notation} (see also \texttt{00\_notation\_register.md}).
The ambient dimension is denoted $n\ge 2$.  We write $|A|=\mathcal L^n(A)$
for the Lebesgue measure of a Borel set $A\subset\R^n$,
$\mathcal H^k$ for the $k$-dimensional Hausdorff measure, and we
identify two sets that differ by an $\mathcal L^n$-null set.

\subsubsection{Finite-perimeter sets and reduced boundary}
\label{ssec:prelim-fp}

\begin{definition}[Finite perimeter and reduced boundary]
\label{def:perimeter}
A Lebesgue-measurable set $E\subset\R^n$ has \emph{finite perimeter}
in an open set $U\subset\R^n$ if the indicator $\mathbf 1_E$ belongs
to $BV(U)$, i.e.\ if
\[
   P(E;U):=\sup\Bigl\{
      \int_E \Div\varphi\,dx :
      \varphi\in C^1_c(U;\R^n),\ \|\varphi\|_\infty\le1
   \Bigr\}<\infty.
\]
When $U=\R^n$ we write $P(E):=P(E;\R^n)$.  The
\emph{De Giorgi reduced boundary} $\partial^*E$ is the set of points
$x\in\R^n$ at which the limit
\[
   \nu_E(x)
   :=\lim_{r\downarrow0}
      \frac{D\mathbf 1_E(B_r(x))}{|D\mathbf 1_E|(B_r(x))}
\]
exists in the unit sphere $\mathbb S^{n-1}$; $\nu_E(x)$ is the
\emph{measure-theoretic outer unit normal} of $E$ at $x$.
By De Giorgi's structure theorem,
\[
   P(E;B)=\mathcal H^{n-1}(\partial^*E\cap B)
   \qquad\text{for every Borel } B\subset\R^n,
\]
and $\partial^*E$ is countably $(n-1)$-rectifiable.
\end{definition}

For the canonical reference see
\cite[Def.~12.1, Thm.~15.9]{Maggi2012}; equivalently
\cite[Def.~3.1, Thm.~3.59]{AFP2000}.  Throughout the manuscript
$P(E)$ always denotes the De Giorgi perimeter and $\nu_E$ (or $\nurho$
when $E=\Erho$) the measure-theoretic outer normal, consistently with
\S 7 of the notation register.

\begin{proposition}[$L^1_{\rm loc}$ lower semicontinuity]
\label{prop:lsc-perimeter}
Let $(E_j)_{j\ge1}$ be measurable subsets of $\R^n$ with
$\mathbf 1_{E_j}\to\mathbf 1_E$ in $L^1_{\rm loc}(\R^n)$ as
$j\to\infty$.  Then for every open $U\subset\R^n$,
\[
   P(E;U)\le\liminf_{j\to\infty}P(E_j;U).
\]
In particular, if $\sup_j P(E_j)<\infty$ and $|E_j|$ is uniformly
bounded, then $(E_j)$ has a subsequence converging in $L^1_{\rm loc}$
to a finite-perimeter set $E$ with $P(E)\le\liminf P(E_j)$.
\end{proposition}

\begin{proof}[Reference]
Lower semicontinuity is
\cite[Prop.~12.15]{Maggi2012} (equivalently
\cite[Thm.~3.7]{AFP2000}); the compactness in $L^1_{\rm loc}$ is
\cite[Thm.~12.26]{Maggi2012}.
\end{proof}

\subsubsection{Coarea, BV coarea, and Sobolev--Sard for the torsion
            function}
\label{ssec:prelim-coarea}

\begin{theorem}[BV coarea formula]\label{thm:bv-coarea}
Let $v\in BV(U)$ on an open set $U\subset\R^n$.  Then $\{v>t\}$ has
finite perimeter in $U$ for $\mathcal L^1$-a.e.\ $t\in\R$, and for
every Borel function $\varphi:U\to[0,\infty]$
\begin{equation}\label{eq:bv-coarea}
   \int_U \varphi\,d|Dv|
   =\int_\R\Bigl(\int_{\partial^*\{v>t\}\cap U}\varphi\,d\mathcal H^{n-1}
      \Bigr)\,dt.
\end{equation}
If $v\in W^{1,1}(U)$ (or $v\in H^1_0(U)$ extended by zero outside
$U$), then $d|Dv|=|\nabla v|\,d\mathcal L^n$ and~\eqref{eq:bv-coarea}
reduces to the Federer--Fleming coarea formula
\begin{equation}\label{eq:fed-fleming}
   \int_U\varphi\,|\nabla v|\,dx
   =\int_\R\Bigl(\int_{\partial^*\{v>t\}\cap U}\varphi\,d\mathcal H^{n-1}
      \Bigr)\,dt.
\end{equation}
\end{theorem}

\begin{proof}[Reference]
This is \cite[Thm.~13.1 and Thm.~18.1]{Maggi2012}; equivalently
\cite[Thm.~3.40 and (3.40)(b)]{AFP2000}.  The vanishing of the
contribution of the critical set $\{\nabla v=0\}$ to the
right-hand side of~\eqref{eq:bv-coarea} is \cite[Thm.~3.40(b)]{AFP2000}.
\end{proof}

\begin{corollary}[Coarea for the torsion function]
\label{cor:coarea-torsion}
Let $\Om\subset\R^n$ be open with $0<|\Om|<\infty$ and let $u=u_\Om$ be
the torsion function (\S 3 of the notation register).  Then $u\in
H^1_0(\Om)\cap L^\infty(\Om)$, and after extending $u$ by zero on
$\R^n\setminus\Om$ we have $u\in BV_{\rm loc}(\R^n)\cap L^\infty(\R^n)$.
For every Borel $\varphi:\R^n\to[0,\infty]$,
\begin{equation}\label{eq:coarea-u}
   \int_{\R^n}\varphi\,|\nabla u|\,dx
   =\int_0^{\|u\|_\infty}
      \Bigl(\int_{\partial^*\{u>t\}}\varphi\,d\mathcal H^{n-1}\Bigr)dt.
\end{equation}
In particular the super-level set $E_t=\{u>t\}$ has finite perimeter
for $\mathcal L^1$-a.e.\ $t\in(0,\|u\|_\infty)$, and the function
$m(t):=|E_t|$ is non-increasing, absolutely continuous on
$(0,\|u\|_\infty)$ in the strict sense made precise in
\S\ref{ssec:level-set-prelim} below.
\end{corollary}

\begin{theorem}[Sobolev--Sard for the torsion function;
                Figalli--Maggi 2011, App.~A, Thm.~A.1]
\label{thm:sobolev-sard}
Let $u=u_\Om$ be as in Corollary~\ref{cor:coarea-torsion}.  There is a
Borel set $\mathcal R\subset(0,\|u\|_\infty)$ of full $\mathcal L^1$
measure such that for every $t\in\mathcal R$:
\begin{enumerate}
   \item[(i)] $E_t=\{u>t\}$ has finite perimeter and
       $m'(t)=-\int_{\partial^*E_t}|\nabla u|^{-1}\,d\mathcal H^{n-1}
        \in(-\infty,0)$;
   \item[(ii)] $|\nabla u|>0$ for $\mathcal H^{n-1}$-a.e.
       $x\in\partial^*E_t$ and the De Giorgi outer normal is
       $\nu_{E_t}(x)=-\nabla u(x)/|\nabla u(x)|$ at $\mathcal H^{n-1}$
       -a.e.\ such $x$;
   \item[(iii)] (Volume-radius parametrisation,
       cf.\ register \S 7.)  The function
       $\rho\mapsto t(\rho)$, characterised by
       $|\{u>t(\rho)\}|=\omega_n\rho^n$, is absolutely continuous on
       $[\rstar,1]$ with $t'(\rho)\le0$, and the normal velocity
       \[
          V_\rho(x):=\frac{-t'(\rho)}{|\nabla u(x)|}
       \]
       is well-defined and positive for $\mathcal H^{n-1}$-a.e.\
       $x\in\partial^*\Erho$, for $\mathcal L^1$-a.e.\
       $\rho\in[\rstar,1]$.
\end{enumerate}
\end{theorem}

\begin{proof}[Reference]
(i)--(ii) are \cite[App.~A, Thm.~A.1]{FigalliMaggi2011}; see also
\cite[Lem.~17.1]{Maggi2012}.  (iii) follows by composition: the
inverse-distribution function $\rho\mapsto t(\rho)$ is the strict-AC
inverse of the strict-AC function $t\mapsto (m(t)/\omega_n)^{1/n}$
(use \cite[Thm.~3.99]{AFP2000} on monotone AC functions), and the
formula for $V_\rho$ then follows from the chain rule and
$m'(t)=-\int|\nabla u|^{-1}$.  The level-set identity draft
\texttt{WIP\_LevelSetIdentity.tex} \S 4 records the same statement.
\end{proof}

\subsubsection{Divergence theorem on sets of finite perimeter}
\label{ssec:div-thm}

\begin{theorem}[Gauss--Green; Maggi Thm.~15.5]
\label{thm:gauss-green}
Let $E\subset\R^n$ have finite perimeter and let
$X\in C^1_c(\R^n;\R^n)$.  Then
\begin{equation}\label{eq:gauss-green-Cc}
   \int_E \Div X\,dx
   =\int_{\partial^*E}X\cdot\nu_E\,d\mathcal H^{n-1}.
\end{equation}
More generally, if $E$ has finite perimeter and finite Lebesgue
measure and $X\in C^1(\R^n;\R^n)\cap L^\infty(E;\R^n)$ satisfies
$\Div X\in L^1(E)$, then~\eqref{eq:gauss-green-Cc} continues to hold
(both sides being finite), as soon as the support of $X$ on
$\partial^*E$ has finite $\mathcal H^{n-1}$ measure (e.g.\ when $E$
is bounded).
\end{theorem}

\begin{proof}[Reference]
The $C^1_c$ statement is \cite[Thm.~15.5]{Maggi2012} or
\cite[Thm.~3.36]{AFP2000}.  The extension to $C^1$ fields bounded on
$E$ with $\Div X\in L^1(E)$ follows by multiplying by a smooth
compactly supported cutoff $\chi_R$, $\chi_R=1$ on $B_R$, sending
$R\to\infty$, and using dominated convergence on both sides; this is
\cite[Thm.~16.3]{Maggi2012} or the form used in
\cite[Ch.~17]{Maggi2012}.
\end{proof}

We will need a slight strengthening allowing $X$ to have a single
locally integrable singularity inside $E$.  The following lemma
is load-bearing for the oriented radial trace estimate of
Agent~11 (\S\ref{thm:radial-trace} below cross-references it).

\begin{lemma}[Truncation lemma for fields with locally integrable
              singularity]\label{lem:truncation}
Let $E\subset\R^n$ have finite perimeter and finite Lebesgue measure,
let $z\in\R^n$, and let $X:\R^n\setminus\{z\}\to\R^n$ be of class
$C^1$ with the following properties:
\begin{itemize}
   \item[(a)] $X\in L^\infty(\R^n;\R^n)$, i.e.\ $|X(x)|\le M_0$ for
       $x\ne z$, for some constant $M_0\ge0$;
   \item[(b)] $\Div X\in L^1_{\rm loc}(\R^n\setminus\{z\})$, and
       there is a constant $K_0\ge0$ such that
       \[
          |\Div X(x)|\le \frac{K_0}{|x-z|}
          \qquad\text{for all }x\ne z;
       \]
   \item[(c)] $E$ is bounded, $E\subset B_R$ for some $R\ge1$.
\end{itemize}
Then $\Div X\in L^1(E)$ (since $n\ge2$) and
\begin{equation}\label{eq:trunc-flux}
   \int_E\Div X\,dx
   =\int_{\partial^*E}X\cdot\nu_E\,d\mathcal H^{n-1},
\end{equation}
where the right-hand side is absolutely integrable.

The prototypical $X$ to which this lemma applies is
\[
   X(x)
   =w(|x-z|)\,\frac{x-z}{|x-z|},
\]
where $w:[0,\infty)\to\R$ is bounded with $w(0)=0$ in the natural
sense (so that $|X|\le\|w\|_\infty$).  In dimension $n\ge2$ such $X$
has $\Div X(x)=w'(|x-z|)+(n-1)w(|x-z|)/|x-z|$ for $x\ne z$, whose
absolute value is $O(|x-z|^{-1})$ if $w$ is bounded with bounded
$w'$.
\end{lemma}

\begin{proof}
We may translate so that $z=0$.  The function $x\mapsto|x|^{-1}$ is
locally integrable in $\R^n$ for $n\ge2$, so (b) gives
$\Div X\in L^1(B_{2R})\supset L^1(E)$.  This proves the first
assertion.

Now choose a fixed cutoff $\eta\in C^\infty(\R)$ with $\eta\equiv0$
on $(-\infty,1]$, $\eta\equiv1$ on $[2,\infty)$ and $0\le\eta\le1$,
and for $\eps\in(0,1)$ set
\[
   \eta_\eps(x):=\eta\bigl(\eps^{-1}|x|\bigr),
   \qquad
   X_\eps(x):=\eta_\eps(x)\,X(x).
\]
Then $X_\eps\equiv 0$ on $B_\eps$, $X_\eps\equiv X$ on $\R^n\setminus
B_{2\eps}$, and $X_\eps\in C^1(\R^n;\R^n)$ (the singularity at $0$
has been removed).  Also $X_\eps$ is bounded by $M_0$ on $\R^n$ and
has compact support if we further multiply by a smooth cutoff
$\chi_R\in C^\infty_c(\R^n)$ with $\chi_R\equiv 1$ on $B_R$,
$\chi_R\equiv0$ on $\R^n\setminus B_{2R}$, $|\nabla\chi_R|\le2/R$;
since $E\subset B_R$, this further cutoff does not affect any of the
integrals over $E$ or $\partial^*E$.

Apply Theorem~\ref{thm:gauss-green} (the $C^1_c$ version) to
$\chi_R X_\eps$:
\begin{equation}\label{eq:trunc-step}
   \int_E \Div(\chi_R X_\eps)\,dx
   =\int_{\partial^*E}(\chi_R X_\eps)\cdot\nu_E\,d\mathcal H^{n-1}.
\end{equation}
Since $\chi_R\equiv1$ on a neighbourhood of $E\subset B_R$, the
left-hand side equals $\int_E\Div X_\eps\,dx$ and the right-hand side
equals $\int_{\partial^*E}X_\eps\cdot\nu_E\,d\mathcal H^{n-1}$.

We now let $\eps\downarrow0$ and verify the two convergences:

\emph{Right-hand side.}  $X_\eps\to X$ pointwise on
$\R^n\setminus\{0\}$, hence $\mathcal H^{n-1}$-a.e.\ on $\partial^*E$
(the single point $0$ has $\mathcal H^{n-1}(\{0\})=0$ since
$n\ge2$).  By (a), $|X_\eps\cdot\nu_E|\le M_0$ pointwise; since
$\mathcal H^{n-1}(\partial^*E)=P(E)<\infty$, dominated convergence
gives
\[
   \int_{\partial^*E}X_\eps\cdot\nu_E\,d\mathcal H^{n-1}
   \longrightarrow
   \int_{\partial^*E}X\cdot\nu_E\,d\mathcal H^{n-1},
\]
and the right-hand integral is absolutely convergent.

\emph{Left-hand side.}  A direct computation gives
\[
   \Div X_\eps(x)
   =\eta_\eps(x)\,\Div X(x)
   +\eps^{-1}\eta'\bigl(\eps^{-1}|x|\bigr)\,
       \frac{x}{|x|}\cdot X(x).
\]
The first summand converges pointwise on $\R^n\setminus\{0\}$ to
$\Div X(x)$ as $\eps\downarrow0$.  It is bounded in absolute value
by $|\Div X(x)|\le K_0/|x|$, which is integrable on $E\subset B_R$.
Dominated convergence gives
\[
   \int_E\eta_\eps\,\Div X\,dx
   \longrightarrow\int_E\Div X\,dx.
\]
For the second summand,
\[
   \Bigl|\eps^{-1}\eta'(\eps^{-1}|x|)\,
         \tfrac{x}{|x|}\cdot X(x)\Bigr|
   \le \frac{\|\eta'\|_\infty M_0}{\eps}\,\mathbf 1_{\{\eps\le|x|\le2\eps\}}.
\]
Hence
\[
   \Bigl|\int_E
      \eps^{-1}\eta'(\eps^{-1}|x|)\,\tfrac{x}{|x|}\cdot X(x)\,dx\Bigr|
   \le
   \frac{\|\eta'\|_\infty M_0}{\eps}\,|B_{2\eps}|
   =\|\eta'\|_\infty M_0\,\omega_n\,2^n\eps^{n-1}.
\]
For $n\ge2$ this tends to $0$ as $\eps\downarrow0$.

Combining, the left-hand side of~\eqref{eq:trunc-step} converges to
$\int_E\Div X\,dx$ and the right-hand side converges to
$\int_{\partial^*E}X\cdot\nu_E\,d\mathcal H^{n-1}$.
This proves~\eqref{eq:trunc-flux}.
\end{proof}

\begin{remark}[Why this matters]\label{rem:trunc-use}
Lemma~\ref{lem:truncation} is used twice in the manuscript:

\noindent(i) (Agent~11, oriented radial trace.) Applied to
\[
   X(x)=w(|x-z|)\,\frac{x-z}{|x-z|},
   \qquad
   w(r)=\Bigl|\tfrac{r}{\rho}-1\Bigr|,
\]
on a bounded torsion level $\Erho\subset B_R$, with the centre
$z=\zbar$ or, more generally, the Fusco--Julin centre on a good
isoperimetric level.  Note that $w$ is bounded by $C(R,\rstar)$ on
$\partial^*\Erho$ (since $|x-z|\le R+R'$ for some
$R'=R'(R,\rstar)$), and a direct computation gives
\[
   \Div X(x)=w'(|x-z|)+(n-1)\frac{w(|x-z|)}{|x-z|},
\]
which is bounded by $C(R,\rstar)(1+|x-z|^{-1})$ pointwise.
Hypotheses (a)--(c) of Lemma~\ref{lem:truncation} are met (with
$M_0=\|w\|_\infty$, $K_0=(n-1)\|w\|_\infty$ and a harmless
$L^\infty$ contribution from $w'$).

\noindent(ii) (Identification of the FJ oscillation index, used in
Agent~4 / Agent~11 background.)  Applied to $X(x)=(x-z)/|x-z|$ on
$\Erho$, hypothesis (b) holds with $K_0=n-1$ since
$\Div((x-z)/|x-z|)=(n-1)/|x-z|$ for $x\ne z$; the divergence
identity
\[
   \int_{\partial^*\Erho}
      \Bigl|\nurho-\frac{x-z}{|x-z|}\Bigr|^2\,d\mathcal H^{n-1}
   =2\,\beta(\Erho,z)^2,
\]
quoted in~\eqref{eq:beta-divergence} below, is a direct consequence
together with the identity $|\xi-\eta|^2=2-2\xi\cdot\eta$ for unit
vectors.
\end{remark}

\subsubsection{Lipschitz truncation flux identity for the torsion
            function}
\label{ssec:flux-torsion}

For the torsion function $u_\Om$ on a regular level $\Erho$, the
boundary trace of $\nabla u$ on $\partial^*\Erho$ is not a classical
trace.  The following identity replaces the formal application of
the divergence theorem to $\nabla u$ on $\Erho$; it is used in
Agent~9 to derive the deficit identity.

\begin{lemma}[Flux $=$ volume; cf.\ Maggi Ch.~17]
\label{lem:flux-volume}
Let $\Om\subset\R^n$ be open with $0<|\Om|<\infty$ and let
$u=u_\Om$.  For $\mathcal L^1$-a.e.\ $t\in(0,\|u\|_\infty)$, the
super-level $E_t=\{u>t\}$ has finite perimeter and
\begin{equation}\label{eq:flux-volume}
   \int_{\partial^*E_t}|\nabla u|\,d\mathcal H^{n-1}=|E_t|.
\end{equation}
\end{lemma}

\begin{proof}
This is proved in
\texttt{WIP\_LevelSetIdentity.tex}~Lem.~3.x via Lipschitz
truncation of the test function $\eta_\eps(x):=\phi_\eps(u(x)-t)$
with $\phi_\eps(s)=\min(1,\max(0,s/\eps))$, testing the weak
equation $-\Delta u=1$ in $H^1_0(\Om)$ against $\eta_\eps$, applying
the coarea formula~\eqref{eq:fed-fleming} to the right-hand side
$\int|\nabla u\cdot\nabla\eta_\eps|=\eps^{-1}\int_{\{t<u<t+\eps\}}
|\nabla u|^2$, and passing to the limit $\eps\downarrow 0$ at every
Lebesgue point of $t\mapsto\int_{\partial^*E_t}|\nabla u|$.  The same
calculation, with $|\nabla u|^2$ replaced by $|\nabla u|^p$ for
suitable $p$, gives the rigorous interpretation of
$\int_{\partial^*E_t}\nabla u\cdot\nu_{E_t}\,d\mathcal H^{n-1}
=-\int_{\partial^*E_t}|\nabla u|$ as the boundary trace of
$\nabla u$ on $\partial^*E_t$ in $L^1$ sense; see
\cite[Ch.~17]{Maggi2012}.
\end{proof}

\subsubsection{Reshetnyak lower semicontinuity and finite-perimeter
            passage of first variations}
\label{ssec:reshetnyak}

\begin{theorem}[Reshetnyak lower semicontinuity for linear-growth
                integrands]\label{thm:reshetnyak}
Let $U\subset\R^n$ be open and let
$f:U\times\R^n\to[0,\infty)$ be continuous, $f(x,\cdot)$ convex and
positively $1$-homogeneous for every $x\in U$.  If
$E_j,E\subset\R^n$ are sets of finite perimeter with
$\mathbf 1_{E_j}\to\mathbf 1_E$ in $L^1_{\rm loc}(U)$, and if the
distributional gradients $D\mathbf 1_{E_j}$ converge weakly-$*$ in
$U$ to $D\mathbf 1_E$, then
\[
   \int_{\partial^*E\cap U}f\bigl(x,\nu_E(x)\bigr)\,d\mathcal H^{n-1}
   \le\liminf_{j\to\infty}\int_{\partial^*E_j\cap U}
        f\bigl(x,\nu_{E_j}(x)\bigr)\,d\mathcal H^{n-1}.
\]
\end{theorem}

\begin{proof}[Reference]
\cite[Thm.~2.38]{AFP2000}.
\end{proof}

This is the workhorse used to pass smooth-flow first-variation
identities to the finite-perimeter setting: the smooth-flow proof of
the metric first-variation bound is performed in
\texttt{WIP\_MetricFramework.tex}, and the four-step approximation
of \texttt{Plan 2/gmt-inputs-for-metric-closure.md} \S 5 turns it
into a finite-perimeter statement.  Concretely the approximation is:

\begin{enumerate}
   \item approximate $u_\Om$ by $u_j\in C^\infty(\Om)$ with $u_j\to
       u$ in $W^{1,2}(\Om)$ and, after a Sobolev--Sard reduction to a
       full-measure set of levels, in the coarea sense on the fixed
       annulus $\rho\in[\rstar,1]$;
   \item form the smooth nested family
       $E_{\rho,j}:=\{u_j>t_j(\rho)\}$, with $t_j(\rho)$ the
       inverse-distribution function of $u_j$;
   \item apply the smooth first-variation bound (proved in
       \texttt{WIP\_MetricFramework.tex} and recalled in
       \texttt{Plan 2/gmt-inputs-for-metric-closure.md} \S 4) to
       $\rho\mapsto E_{\rho,j}$;
   \item pass to the limit using Proposition~\ref{prop:lsc-perimeter}
       (lower semicontinuity of perimeter),
       Theorem~\ref{thm:reshetnyak} (Reshetnyak lower semicontinuity
       for the normal integrals
       $\int|V_\rho-H_{\zbar,\rho}-a\cdot\nurho|\,d\mathcal H^{n-1}$),
       and the standard lower semicontinuity of metric derivatives
       for $L^1$-convergent curves \cite[Thm.~2.2.1]{AGS2008}.
\end{enumerate}

The detailed passage is contained in the metric-framework section
(Agent~10).  We isolate the following corollary, which is the form
needed by Agents 11--12.

\begin{corollary}[Reshetnyak passage for the homothetic defect]
\label{cor:reshetnyak-homot}
Let $E_j,E\subset B_R$ have finite perimeter with $|E_j|=|E|=
\omega_n\rho^n$ and $\mathbf 1_{E_j}\to\mathbf 1_E$ in $L^1(\R^n)$,
together with $P(E_j)\to P(E)$.  Let $z_j\to z\in\R^n$ be a sequence
of centres.  For every $a\in\R^n$,
\[
   \int_{\partial^*E}
      \bigl|V-H_{z,\rho}-a\cdot\nu_E\bigr|\,d\mathcal H^{n-1}
   \le\liminf_{j\to\infty}
   \int_{\partial^*E_j}
      \bigl|V_j-H_{z_j,\rho}-a\cdot\nu_{E_j}\bigr|\,d\mathcal H^{n-1},
\]
whenever $V_j\to V$ in $L^1(\mathcal H^{n-1}\llcorner\partial^*E_j)$ in
the sense of weak convergence of varifolds.
\end{corollary}

\begin{proof}[Sketch]
Apply Theorem~\ref{thm:reshetnyak} to the integrand
$f(x,\nu)=|V(x)-H_{z,\rho}(x)-a\cdot\nu|$, which is continuous in
$x$, convex in $\nu$, and positively $1$-homogeneous (because
$H_{z,\rho}(x)=(x-z)\cdot\nu(x)/\rho$ is linear in $\nu$).  The
convergence $P(E_j)\to P(E)$ together with $\mathbf 1_{E_j}\to
\mathbf 1_E$ in $L^1$ gives, by~\cite[Thm.~12.30]{Maggi2012}, weak-$*$
convergence of $D\mathbf 1_{E_j}$ to $D\mathbf 1_E$ in the sense
needed by Theorem~\ref{thm:reshetnyak}.  Convergence of $V_j$ in
varifold sense is the standing hypothesis.
\end{proof}

\subsubsection{Measurable selection of levelwise centres}
\label{ssec:meas-sel}

The Fusco--Julin centre package (Agent~11) attaches to every level
$\rho\in[\rstar,\rho_\delta]$ in the good set $G_I\subset
[\rstar,\rho_\delta]$ at least one point $z\in\R^n$ such that the
quantitative isoperimetric inequality is realised at $z$; this point
is unique only up to a translation in the bad set, and to use it
inside metric-derivative integrals one must choose it in a Borel way.

\begin{lemma}[Borel selection of centres]
\label{lem:borel-selection}
Let $G\subset[\rstar,1]$ be Borel and for each $\rho\in G$ let
$Z(\rho)\subset\R^n$ be a non-empty closed set; suppose the
graph $\{(\rho,z)\in G\times\R^n:z\in Z(\rho)\}$ is a Borel subset of
$G\times\R^n$ and $\sup_{\rho\in G}\sup_{z\in Z(\rho)}|z|<\infty$.
Then there exists a bounded Borel map $\sigma:G\to\R^n$ with
$\sigma(\rho)\in Z(\rho)$ for every $\rho\in G$.
\end{lemma}

\begin{proof}
The graph of $Z$ is a Borel subset of the standard Borel space
$G\times\R^n$ with non-empty closed $\rho$-sections.  The
Kuratowski--Ryll-Nardzewski measurable selection theorem
\cite[Thm.~III.6 and Thm.~III.8]{CastaingValadier} (also Federer,
\cite[\S 2.13]{Federer1969}) provides a Borel selection
$\sigma:G\to\R^n$.  Boundedness of $\sigma$ follows from the standing
hypothesis that the union $\bigcup_{\rho\in G}Z(\rho)$ is bounded.

Outside $G$ (i.e.\ in the bad set), set $\sigma(\rho):=0$; this
extends $\sigma$ to a bounded Borel map on $[\rstar,1]$.
\end{proof}

\begin{remark}[Application]\label{rem:borel-application}
In the manuscript, $Z(\rho)$ is taken to be the set of
$\eta_I$-almost minimisers in the scaled Fusco--Julin inequality
applied to $\Erho$: explicitly, the set of $z\in\R^n$ such that
\[
   |\Erho\Delta B_\rho(z)|^2+\rho^{n+1}\beta(\Erho,z)^2
   \le 2\,C_n\,\rho^{n+1}\bigl(P(\Erho)-P(B_\rho)\bigr),
\]
with $C_n$ the Fusco--Julin constant
\cite[Thm.~1.1]{FJ2014}; this is non-empty by Fusco--Julin
itself, closed by lower semicontinuity of the integrals in $z$
(continuity of $z\mapsto|\Erho\Delta B_\rho(z)|$ is elementary, and
continuity of $z\mapsto\beta(\Erho,z)^2=P(\Erho)-(n-1)\int_{\Erho}
|x-z|^{-1}dx$ follows from dominated convergence using the bound
$|x-z|^{-1}\le|x-z_0|^{-1}+|z-z_0||x-z|^{-1}|x-z_0|^{-1}$), and
bounded in $z$ because the Fraenkel part forces $|z|\le R+\rho$ on
$E\subset B_R$.  Lemma~\ref{lem:borel-selection} then provides a
bounded Borel field $\rho\mapsto z_\rho$ on the good set $G_I$.
The Borel-measurability of the graph follows from the fact that
both $(z,E)\mapsto|E\Delta B_\rho(z)|$ and $(z,E)\mapsto
\beta(E,z)^2$ are Borel maps in $z$ for fixed $E$ (in fact
continuous), so the level set is closed and the graph is therefore
Borel.
\end{remark}

\subsubsection{Two auxiliary identities used downstream}
\label{ssec:aux-identities}

For convenience we record two identities used by Agents~11 and~12.

\begin{lemma}[FJ oscillation index identity]
\label{lem:beta-divergence}
Let $E\subset\R^n$ have finite perimeter and finite measure with
$z\notin\R^n$ playing the role of a centre.  Define
\[
   \beta(E,z)^2:=P(E)-(n-1)\int_E|x-z|^{-1}\,dx.
\]
Then $\beta(E,z)^2\ge 0$ and
\begin{equation}\label{eq:beta-divergence}
   \int_{\partial^*E}
      \Bigl|\nu_E(x)-\frac{x-z}{|x-z|}\Bigr|^2\,d\mathcal H^{n-1}
   =2\,\beta(E,z)^2.
\end{equation}
\end{lemma}

\begin{proof}
The identity $|\nu_E(x)-e_r|^2=2-2e_r\cdot\nu_E(x)$ for unit vectors
$e_r$ gives
\[
   \int_{\partial^*E}|\nu_E-e_r|^2
   =2P(E)-2\int_{\partial^*E}e_r\cdot\nu_E\,d\mathcal H^{n-1}.
\]
Apply Lemma~\ref{lem:truncation} (with $z$ in place of the origin) to
the vector field $X(x):=(x-z)/|x-z|$.  Hypothesis (a) holds with
$M_0=1$, and a direct computation gives $\Div X(x)=(n-1)/|x-z|$ for
$x\ne z$, so (b) holds with $K_0=n-1$.  (c) holds because $E$ has
finite measure and, in the manuscript application, $E\subset B_R$.
Therefore
\[
   \int_{\partial^*E}\frac{x-z}{|x-z|}\cdot\nu_E\,d\mathcal H^{n-1}
   =\int_E\Div X\,dx
   =(n-1)\int_E|x-z|^{-1}\,dx.
\]
Combining gives~\eqref{eq:beta-divergence}.  Non-negativity is then
immediate from the left-hand side, and equality with the right-hand
side proves $\beta(E,z)^2\ge0$.
\end{proof}

\begin{lemma}[Total flux of $H_{z,\rho}$]\label{lem:H-total-flux}
For $z\in\R^n$ and a set $\Erho$ of finite perimeter with
$|\Erho|=\omega_n\rho^n$,
\[
   \int_{\partial^*\Erho}H_{z,\rho}\,d\mathcal H^{n-1}
   =\frac{n|\Erho|}{\rho}=n\omega_n\rho^{n-1}=P(B_\rho).
\]
\end{lemma}

\begin{proof}
Apply Theorem~\ref{thm:gauss-green} to $X(x):=(x-z)/\rho\in
C^1(\R^n;\R^n)$, which is $C^1$ everywhere and bounded on $\Erho$
(any finite-measure set is $\mathcal H^{n-1}$-essentially bounded on
its reduced boundary; in the manuscript $\Erho\subset B_R$).
$\Div X=n/\rho$, so the divergence theorem gives
$\int_{\partial^*\Erho}X\cdot\nurho\,d\mathcal H^{n-1}=n|\Erho|/\rho$.
Since $(x-z)\cdot\nurho(x)/\rho=H_{z,\rho}(x)$, the lemma follows.
\end{proof}

\subsubsection{Summary of the GMT preliminaries and how they are used}
\label{ssec:summary}

For the convenience of the reader we summarise which GMT facts enter
each Plan~2 section.

\begin{itemize}
   \item \emph{Level-set deficit identity (Agent~9).}
       BV coarea (Theorem~\ref{thm:bv-coarea}), the Sobolev--Sard
       statement for the torsion function
       (Theorem~\ref{thm:sobolev-sard}), and the Lipschitz-truncation
       flux identity (Lemma~\ref{lem:flux-volume}).

   \item \emph{Metric framework (Agent~10).}
       Lower semicontinuity of perimeter
       (Proposition~\ref{prop:lsc-perimeter}), the divergence
       theorem (Theorem~\ref{thm:gauss-green}), and Reshetnyak lower
       semicontinuity together with the smooth-to-finite-perimeter
       approximation (Theorem~\ref{thm:reshetnyak},
       Corollary~\ref{cor:reshetnyak-homot}).

   \item \emph{Fusco--Julin centre and oriented radial trace
       (Agent~11).}
       Truncation lemma (Lemma~\ref{lem:truncation}); the FJ
       oscillation-index identity (Lemma~\ref{lem:beta-divergence});
       and the Borel selection
       (Lemma~\ref{lem:borel-selection}).

   \item \emph{Integrated action and weighted metric trace
       (Agent~12).}
       The same divergence and BV coarea tools, plus the absolute
       continuity of $\rho\mapsto[\Frho]$ proved in the metric
       framework section.
\end{itemize}

All boundary integrals in the manuscript are on De Giorgi reduced
boundaries; all centre fields are Borel via
Lemma~\ref{lem:borel-selection}; the truncation argument of
Lemma~\ref{lem:truncation} is the rigorous form of the
``$x/|x|$ argument'' invoked at several places in Plan~2.

% End of Preliminaries -- Geometric Measure Theory (Agent 2).
